% Build in this directory: pdflatex -interaction=nonstopmode -halt-on-error class5_handout.tex
\documentclass[a4paper,10pt,pdflatex,ja=standard,jafont=ipaex-type1,
  magstyle=real,scale=1]{bxjsarticle}
\usepackage{lmodern}
\usepackage{amsmath,amssymb,mathtools}
\usepackage{geometry}
\geometry{reset,left=22mm,right=22mm,top=20mm,bottom=17mm,
  headheight=13pt,headsep=6mm,footskip=10mm}
\usepackage{xcolor}
\usepackage{enumitem}
\usepackage{array,booktabs,tabularx}
\usepackage{fancyhdr}

\definecolor{ink}{HTML}{193A5C}
\definecolor{accent}{HTML}{237D82}
\definecolor{soft}{HTML}{EAF3F3}
\definecolor{muted}{HTML}{52616B}
\newcommand{\coursename}{解析学 II}
\newcommand{\handoutnumber}{5}
\pagestyle{fancy}
\fancyhf{}
\fancyhead[L]{\small\color{ink}\coursename{}・授業要約と演習}
\fancyhead[R]{\small\color{ink}第\handoutnumber 回}
\fancyfoot[C]{\small\color{ink}\thepage/2}
\renewcommand{\headrulewidth}{0.3pt}
\setlength{\parindent}{0pt}
\setlength{\parskip}{2pt}
\setlength{\abovedisplayskip}{5pt}
\setlength{\belowdisplayskip}{5pt}
\setlength{\abovedisplayshortskip}{3pt}
\setlength{\belowdisplayshortskip}{4pt}
\setlist[itemize]{leftmargin=1.8em,itemsep=1pt,topsep=2pt,parsep=0pt}

\newcommand{\handouttitle}[1]{%
  \begin{center}
    {\Large\color{ink}#1}\par\vspace{3mm}
    {\color{accent}\rule{0.88\linewidth}{0.7pt}}
  \end{center}\vspace{-2mm}}
\newcommand{\handout}[2]{%
  \renewcommand{\handoutnumber}{#1}%
  \handouttitle{第#1回\quad #2}}
\newcommand{\topic}[1]{\par\vspace{5pt}%
  {\large\sffamily\mdseries\color{accent}#1}\par\vspace{2pt}}
\newcommand{\keybox}[1]{\par\vspace{3pt}\begingroup
  \setlength{\fboxsep}{4pt}%
  \colorbox{soft}{\parbox{\dimexpr\linewidth-2\fboxsep\relax}{\small #1}}%
  \endgroup\par\vspace{2pt}}
\newenvironment{problems}
  {\begin{enumerate}[leftmargin=2em,itemsep=3pt,topsep=3pt,parsep=0pt,
    font=\color{accent}]}
  {\end{enumerate}}


\newcommand{\examplelabel}[1]{{\sffamily\mdseries\color{ink}#1}}
\newcommand{\answerroom}{\par\vspace{16mm}}

\begin{document}
\fontsize{10pt}{13.5pt}\selectfont
\setlength{\abovedisplayskip}{4pt}
\setlength{\belowdisplayskip}{4pt}
\handout{5}{広義積分}
\topic{特異点や無限遠を極限で扱う}
$f$ が $[a,b)$ の各閉部分区間で可積分であるとき
$$
 \int_a^bf(x)\,dx:=\lim_{t\to b-}\int_a^tf(x)\,dx,
 \qquad
 \int_a^\infty f(x)\,dx:=\lim_{R\to\infty}\int_a^Rf(x)\,dx.
$$
有限の極限が存在するとき「収束する」という。内部の特異点 $c\in(a,b)$ は左右に分け、
両方の広義積分がそれぞれ収束する場合だけ和を取る。
$$
 \int_a^bf=\lim_{t\to c-}\int_a^tf+\lim_{s\to c+}\int_s^bf.
$$
同じ幅を除いて左右を相殺する極限は、通常の広義積分と区別する。

\topic{べき関数の収束判定と比較}
$p\in\mathbb R$ に対して、特異点付近と無限遠では条件が逆になる。
$$
 \int_0^1x^{-p}\,dx=
 \begin{cases}\dfrac1{1-p},&p<1,\\+\infty,&p\ge1,\end{cases}
 \qquad
 \int_1^\infty x^{-p}\,dx=
 \begin{cases}\dfrac1{p-1},&p>1,\\+\infty,&p\le1.\end{cases}
$$
$0\le f\le g$ なら、$\int g<\infty$ は $\int f<\infty$ を導き、
$\int f=+\infty$ は $\int g=+\infty$ を導く。
符号が変わる関数では、まず $\int|f|<\infty$（絶対収束）を調べる。

\topic{例：端点で発散しても積分は収束する}
$$
 \int_0^1\frac{dx}{\sqrt{1-x}}
 =\lim_{\varepsilon\downarrow0}
 \left[-2\sqrt{1-x}\right]_0^{1-\varepsilon}
 =\lim_{\varepsilon\downarrow0}(2-2\sqrt\varepsilon)=2.
$$
関数の値が無限大になることと、積分が発散することは同じではない。

\topic{非負関数の広義重積分}
特異点を除いた領域を $D$ とし、$D_1\subset D_2\subset\cdots$ を
有界な積分可能領域、$\bigcup_nD_n=D$ とする。各 $D_n$ 上で連続な $f\ge0$ について
$$
 \iint_Df\,dx\,dy:=\lim_{n\to\infty}\iint_{D_n}f\,dx\,dy\in[0,+\infty].
$$
非負の場合、この値は領域を増やす取り方によらない。
符号が変わる場合は $\iint_D|f|<\infty$ を仮定すれば、領域の取り方や累次積分の順序に依存しない。

\topic{例：2変数の特異点を避ける}
$D=(0,1]^2$、$D_\varepsilon=[\varepsilon,1]^2$ とすると
$$
 \iint_D\frac{dx\,dy}{\sqrt{xy}}
 =\lim_{\varepsilon\downarrow0}
 \left(\int_\varepsilon^1x^{-1/2}\,dx\right)^2
 =\lim_{\varepsilon\downarrow0}4(1-\sqrt\varepsilon)^2=4.
$$
\keybox{特異点を列挙し、各部分の極限を別々に調べる。多変数では、非負性または絶対可積分性を確認する。}

\newpage
\handouttitle{第5回\quad 演習問題}
計算過程を記し、必要な条件・積分領域・向きを明示すること。
\begin{problems}
\item $p\in\mathbb R$ に対し、$\int_0^1x^{-p}\,dx$ と $\int_1^\infty x^{-p}\,dx$ の収束条件を、原始関数と極限から導け。$p=1$ も扱うこと。\answerroom
\item $\displaystyle\int_0^1\frac{dx}{\sqrt{x}}$ と $\displaystyle\int_0^1\frac{dx}{1-x}$ の収束・発散を調べ、収束する場合は値を求めよ。\answerroom
\item $\displaystyle\int_1^\infty\frac{dx}{x(x+1)}$ が収束することを示し、その値を求めよ。\answerroom
\item $\displaystyle\int_{-1}^1\frac{dx}{x}$ は通常の広義積分として収束するか。$\lim_{\varepsilon\downarrow0}(\int_{-1}^{-\varepsilon}\!dx/x+\int_\varepsilon^1\!dx/x)$ と比較せよ。\answerroom
\item $p,q\in\mathbb R$ とする。$\displaystyle\iint_{(0,1]^2}x^{-p}y^{-q}\,dx\,dy$ の収束条件と、収束する場合の値を求めよ。\answerroom
\item $\displaystyle\iint_{\mathbb R^2}\frac{dx\,dy}{(1+x^2)(1+y^2)}$ の収束と値を調べよ。累次積分を使える理由も述べよ。\answerroom
\end{problems}
\end{document}
